% PATCH: R2-01 (19/13, 3rd-pass cross-check)
% Reviewer: 2 (redundancy reduction), touching R3-03's (Deivi's) content
% Target section: sections/methodology.tex
% Requirement: Condense repetitive text -- this time a cross-section, cross-author case found
%   while auditing recently-synced content for the redundancy sweep. R3-03's new
%   "Low-Torque Actuation Compensation Strategies" subsection (ssec:ServoCompensation, added by
%   Deivi, synced-to-overleaf) independently re-argued the same claim already made -- with
%   quantitative backing -- in results.tex's "Quantitative Stability Analysis of Locomotion Mode
%   Transitions" subsubsection (sssec:transition_stability, Table~tab:crawl_stability, N=268
%   samples): that the crawl gait is statically/quasi-statically stable by construction because
%   only one leg is ever airborne, keeping the CoM within the support polygon/triangle. Neither
%   author was aware of the other's phrasing at the time of writing.
% Fix: kept the full quantitative claim in results.tex untouched; condensed ServoCompensation's
%   paragraph to state the fact once (needed for its own compensation-strategy argument) plus a
%   cross-reference to the quantitative confirmation, instead of re-deriving the stability
%   argument from scratch. The 3 core compensation techniques and the pending IMU question
%   (R3-03's own scope) are unchanged.
% Status: applied

% Before:
%   "First, gait stability is addressed at the scheduling level rather than through active
%   balance control. The crawl gait described in Section~\ref{ssec:Gait_Decision_Module}
%   advances exactly one leg at a time, following a fixed diagonal sequence (front-left ->
%   rear-right -> front-right -> rear-left); the remaining three legs stay grounded and fixed
%   throughout each swing phase. This is a structural, not a runtime-computed, guarantee:
%   because only one leg is ever airborne, the robot's center of mass is always contained
%   within the support triangle formed by the three stance legs, independently of any per-step
%   positioning error the open-loop actuation may introduce. The gait is therefore
%   quasi-statically stable by construction, which bounds the destabilizing consequence of an
%   individual joint's torque deficit to a single, momentarily unsupported limb rather than to
%   the whole platform."
% After:
\revblue{First, gait stability is addressed at the scheduling level rather than through active balance control. The crawl gait described in Section~\ref{ssec:Gait_Decision_Module} advances exactly one leg at a time, following a fixed diagonal sequence (front-left~$\rightarrow$~rear-right~$\rightarrow$~front-right~$\rightarrow$~rear-left); the remaining three legs stay grounded and fixed throughout each swing phase. This is a structural, not a runtime-computed, guarantee, quantitatively confirmed in Section~\ref{sssec:transition_stability} (Table~\ref{tab:crawl_stability}): because only one leg is ever airborne, this bounds the destabilizing consequence of an individual joint's torque deficit to a single, momentarily unsupported limb rather than to the whole platform.}
